% GENERATED FILE. Edit canonical lessons, then run npm run generate:books.
% canonical-source-hash: fnv1a64:dac8def569fd0a19
% canonical-lessons: PT-C264-ida-e-volta, PT-C264-perdi-o-comboio, PT-C264-revisao-estacao

\chapter{Ida e Volta --- The Ticket, and the Missed Train}
\label{ch:pt-the-ticket-and-the-missed-train}
\hlchaptermodality{264}

\begin{chapteropening}
\textbf{By the end of this chapter:} \emph{I can buy a single or return ticket, and report a problem at the station with perdi o comboio.}

Complete PT-C264-revisao-estacao: find the station, buy the ticket, miss the train and get moving again.
\end{chapteropening}

\section[ida e volta, só de ida]{ida e volta — the going and the coming back}

\label{lesson:PT-C264-ida-e-volta}



\begin{quote}
Say \textquotedblleft{}a ticket\textquotedblright{}. (\emph{Um bilhete.}) Say \textquotedblleft{}I would like\textquotedblright{}. (\emph{Gostaria.})

Now you are standing at a ticket window.
\end{quote}

\subsection*{You'll want to know: ida e volta, só de ida}
\textbf{ida e volta} — \textquotedblleft{}there and back\textquotedblright{}: a return ticket. \textbf{só de ida} — \textquotedblleft{}one way\textquotedblright{}, the going only.

\begin{quote}
\emph{Um bilhete de ida e volta.} — A return ticket. \emph{Só de ida.} — One way only.
\end{quote}

\textbf{a ida} is \textquotedblleft{}the going\textquotedblright{}, from \textbf{ir}; \textbf{a volta} is \textquotedblleft{}the coming back\textquotedblright{}, from your \textbf{voltar}. \textbf{só} means \textquotedblleft{}only\textquotedblright{}. Brazil says \emph{passagem} where Portugal says \emph{bilhete}, and \emph{trem} for \emph{comboio}; \emph{ida e volta} is the same on both sides of the ocean.

\begin{grammarlens}[title={Grammar Lens: queria um bilhete para…}]
The whole request takes one frame:

\begin{quote}
\textbf{Queria um bilhete para} Lisboa, \textbf{de ida e volta}. — I'd like a return ticket to Lisbon.
\end{quote}

\textbf{queria} is the imperfeito of \textbf{querer}, and at a counter in Portugal it is the everyday polite request: \textquotedblleft{}I was wanting\textquotedblright{} softens into \textquotedblleft{}I'd like\textquotedblright{}. \emph{Gostaria} is more formal. Most cities take no article, \emph{para Lisboa}, but a few keep one: \emph{para o Porto}.

And what comes back at you:

\begin{quote}
— \emph{Ida e volta?} — \emph{Sim, ida e volta.} / \emph{Não, só de ida.} — \emph{A que horas parte o comboio?} — \emph{Às dez.}
\end{quote}

\emph{Às} is \emph{a} plus \emph{as}, \textquotedblleft{}at the\textquotedblright{}, before the hour.
\end{grammarlens}

\begin{cousinweb}
\textbf{ir} is inherited from Latin \emph{īre}, \textquotedblleft{}to go\textquotedblright{} (EtymDB 2.1), and \textbf{ida} is built on its participle \emph{ido}, made a feminine noun, \textquotedblleft{}the going\textquotedblright{}; EtymDB 2.1 has no row for the noun itself. \textbf{volta} comes from \textbf{voltar}, inherited from spoken Latin \emph{voltāre}, \textquotedblleft{}to turn round\textquotedblright{} (EtymDB 2.1), from Latin \emph{volūtāre}, \textquotedblleft{}to roll about\textquotedblright{} (Lewis \& Short, under \emph{voluto}). A return is a turning round. \textbf{só} is inherited from Latin \emph{sōlus}, \textquotedblleft{}alone\textquotedblright{} (EtymDB 2.1).
\end{cousinweb}

\subsection*{Your turn}
Say these aloud:

\begin{itemize}
\raggedright
  \item \emph{ida e volta}, then \emph{só de ida}
  \item \emph{Queria um bilhete para Lisboa, de ida e volta.}
  \item the one-way version — \emph{Queria um bilhete para o Porto, só de ida.}
  \item ask the time — \emph{Desculpe, a que horas parte o comboio?}
\end{itemize}

\begin{tcolorbox}[breakable,colback=teal!4,colframe=teal!35!black,title={Before you move on}]
What is a one-way ticket? (\emph{\textbf{Só de ida.}}) Which past tense is \emph{queria}, and why use it at a counter? (The \textbf{imperfeito}: it softens \textquotedblleft{}I want\textquotedblright{} into \textquotedblleft{}I'd like\textquotedblright{}.) Which city keeps its article after \emph{para}? (\emph{\textbf{O Porto}}.)
\end{tcolorbox}
\section[perdi o comboio]{perdi o comboio — when it goes wrong}

\label{lesson:PT-C264-perdi-o-comboio}



\begin{quote}
Ask for a one-way ticket to Porto. (\emph{Queria um bilhete para o Porto, só de ida.})

Now the ticket is in your pocket, and the train has gone without you.
\end{quote}

\subsection*{You'll want to know: perdi o comboio}
\textbf{perdi o comboio} — \textquotedblleft{}I've missed the train\textquotedblright{}.

You have \textbf{perder}, \textquotedblleft{}to lose\textquotedblright{}. Portuguese uses the same verb for a train you missed: you \textbf{lost} it. Its event past is the \emph{-i} of \emph{dormi}: \textbf{perdi}, and for \textquotedblleft{}he\textquotedblright{} or \textquotedblleft{}she\textquotedblright{}, \textbf{perdeu}.

\begin{quote}
\textbf{Perdi o comboio.} — I've missed the train. \textbf{Perdi o autocarro.} — I've missed the bus. \textbf{Perdi a mala.} — I've lost my suitcase.
\end{quote}

Portuguese says \emph{a} mala where English says \emph{my} suitcase: whose it is goes without saying. One warning about the present: \textquotedblleft{}I lose\textquotedblright{} is irregular, \textbf{perco}.

And listen to the English. \textquotedblleft{}I've missed\textquotedblright{} sounds like \emph{tenho perdido}, but chapter fifteen showed what that form means in Portuguese: something that keeps happening lately. \emph{Tenho perdido o comboio} is \textquotedblleft{}I keep missing the train\textquotedblright{}. For one missed train, the plain \textbf{perdi}.

\begin{grammarlens}[title={Grammar Lens: reporting a problem}]
A complaint to a stranger has three parts, and the first is a courtesy:

\begin{quote}
\textbf{Desculpe, tenho um problema.} — Excuse me, I have a problem. \textbf{Perdi o comboio.} — What went wrong. \textbf{Há outro comboio hoje? A que horas?} — What you need.
\end{quote}

Or the fault is not yours: \emph{O comboio não chegou}, \textquotedblleft{}the train didn't come\textquotedblright{}. Say what happened, in the perfeito, then ask the question that gets you moving again.
\end{grammarlens}

\begin{cousinweb}
\textbf{perder} is inherited from Latin \emph{perdere} (EtymDB 2.1), built on \emph{per-}, \textquotedblleft{}through\textquotedblright{}, and \emph{dare}, \textquotedblleft{}to give\textquotedblright{}: to give something away for good, to squander, to lose (Lewis \& Short, under \emph{perdo}). The \emph{dare} inside it is the one that gave \textbf{dar}, \textquotedblleft{}to give\textquotedblright{}.
\end{cousinweb}

\subsection*{Your turn}
Say these aloud:

\begin{itemize}
\raggedright
  \item \emph{Perdi o comboio.} — then \emph{Perdi o autocarro.}
  \item the three parts — \emph{Desculpe, tenho um problema. Perdi o comboio. Há outro comboio hoje?}
  \item the fault that is not yours — \emph{O comboio não chegou.}
  \item a habit, then one time — \emph{Perco sempre o autocarro.} then \emph{Ontem perdi o autocarro.}
\end{itemize}

\begin{tcolorbox}[breakable,colback=teal!4,colframe=teal!35!black,title={Before you move on}]
Say \textquotedblleft{}I've missed the train\textquotedblright{}. (\emph{\textbf{Perdi o comboio.}}) What does \emph{tenho perdido o comboio} say instead? (That it \textbf{keeps happening}.) Name the three parts of a complaint. (A \textbf{courtesy}, \textbf{what went wrong}, \textbf{what you need}.) And what are the two Latin pieces of \emph{perder}? (\emph{\textbf{Per-}} and \emph{\textbf{dare}}, \textquotedblleft{}to give\textquotedblright{}.)
\end{tcolorbox}
\section[(review)]{Review — one morning at the station}

\label{lesson:PT-C264-revisao-estacao}



\begin{quote}
Say \textquotedblleft{}I've missed the train\textquotedblright{}. (\emph{Perdi o comboio.})
\end{quote}

\begin{grammarlens}[title={What you've built}]
Two chapters of travel, and they make one morning. You find the station, you buy the ticket, something goes wrong, and you ask your way out of it.

Every step is a sentence you have said: \emph{Desculpe, como se vai para a estação?} · \emph{Queria um bilhete para Lisboa, de ida e volta.} · \emph{Desculpe, tenho um problema: perdi o comboio.} · \emph{Há outro comboio hoje?}

And when you tell someone about it afterwards, the morning becomes a story: \emph{Havia muitas pessoas na estação, por isso perdi o comboio.} Scene, arrow, event.
\end{grammarlens}

\subsection*{Your turn}
Say these aloud:

\begin{itemize}
\raggedright
  \item find it — \emph{Desculpe, como se vai para a estação?}
  \item the answer, said back — \emph{Em frente, depois à esquerda?}
  \item buy it — \emph{Queria um bilhete para Lisboa, de ida e volta.}
  \item the problem — \emph{Desculpe, tenho um problema. Perdi o comboio.}
  \item tell it afterwards — \emph{Perdi o comboio porque havia muitas pessoas na estação.}
\end{itemize}

\begin{tcolorbox}[breakable,colback=teal!4,colframe=teal!35!black,title={Before you move on}]
What is a return ticket? (\emph{\textbf{Ida e volta.}}) What do you say before the problem itself? (\emph{\textbf{Desculpe, tenho um problema.}}) In \emph{perdi o comboio porque havia muitas pessoas}, which verb is the scene? (\emph{\textbf{Havia}}.)
\end{tcolorbox}
